% 第 12 章 Working With the GUI
\bichapter{使用 GUI}{Working With the GUI}
\label{chap:gui}

Design Vision 是 Synopsys 逻辑综合环境的图形用户界面（GUI）。
它提供在通用技术（GTECH）级与门级查看并分析设计的分析工具，并具备 DC Explorer 工具的全部综合能力。

GUI 提供多种可视化与分析功能，可用于可视化设计数据并分析结果。
GUI 窗口为最常用的综合功能提供菜单命令与对话框；此外，也可在 GUI 或 shell 的命令行中输入任意 \cmd{de\_shell} 命令。

在使用 GUI 分析与排查设计之前，应先熟悉其操作及所提供的各类工具。
下列各节给出首次使用 GUI 前需要了解的一般与具体信息：
\begin{itemize}
  \item 功能与收益
  \item 使用 GUI 窗口
  \item 执行 Floorplan 探索
  \item 使用 IC Compiler II 进行 Floorplan 探索
  \item 分析 RTL
  \item RTL 交叉探测（RTL Cross-Probing）
\end{itemize}

% ============================================================
\section{功能与收益}
\subsection*{Features and Benefits}

Design Vision GUI 提供下列功能：
\begin{itemize}
  \item 面向 Synopsys 逻辑综合的窗口与菜单驱动界面
  \item 分析可视化能力，包括：
  \begin{itemize}
    \item hierarchy browser，用于在设计层次中导航并探索设计结构，包括查看 area attribute
    \item 直方图（histogram），用于可视化各类度量的趋势，例如 slack 与电容
    \item 时序状态摘要与路径检查器（path inspector），用于详细时序路径分析
    \item 路径原理图（path schematic），用于直观检查时序路径、扇入/扇出逻辑以及网连接
    \item 原理图视图（schematic view），用于直观检查高层次与低层次设计连接
    \item 属性查看器（properties viewer），用于查看对象属性并编辑属性值
    \item 版图视图（layout view），用于查看并分析物理设计中的 floorplan 元素
  \end{itemize}
  \item 报告功能，可将报告中的对象与图形视图关联
  \item 集成于 GUI 的 command-line interface，支持以 Design Compiler tool command language（dctcl）的全部命令编写脚本
  \item 基于 HTML 的文档浏览器，用于查看 man 页与联机帮助页
\end{itemize}

借助 Design Vision GUI 的功能，您可以：
\begin{itemize}
  \item 执行 DC Explorer 综合与报告
  \item 探索设计结构
  \item 用 Design Vision 图形分析可视化整体时序
  \item 对正在综合的 block 执行 timing analysis
  \item 对时序路径与相连逻辑做详细可视化分析
  \item 在物理设计版图中直观检查并调试时序路径及某些 floorplan 约束
\end{itemize}

% ============================================================
\section{使用 GUI 窗口}
\subsection*{Using GUI Windows}

GUI 提供灵活的工作环境，包含多个用于不同任务的窗口。
这些 GUI 窗口在您的 X Windows 环境中彼此独立运行，但共享内存中的同一设计、相同的当前时序信息，以及工具中的全局选择数据。

Design Vision 提供下列 GUI 窗口：
\begin{itemize}
  \item Design Vision 主窗口：从 \cmd{de\_shell} 打开 GUI 时自动出现。\\
        默认在 Design Vision 主窗口中显示 hierarchy browser（logic hierarchy view）与 command console（console panel）。
  \item 版图窗口（layout window）：用于可视化设计的物理方面。\\
        要打开版图窗口，选择 \textbf{Window $>$ New Layout Window}。
\end{itemize}

每个 GUI 窗口顶部包含标题栏、菜单栏与若干工具栏，底部为状态栏。
工具栏与状态栏之间的工作区显示视图窗口与面板。
视图窗口提供设计信息的图形或文本视图；面板提供用于设置选项或执行常用任务的交互工具。

可用 Linux 桌面的窗口管理工具调整、移动、最小化或最大化各 GUI 窗口。更多细节见 Design Vision Help。
也可打开同一 GUI 窗口的多个实例，以便并排比较视图或同一视图中的不同设计信息。窗口在整个会话中按序编号。

图~\ref{fig:dv-main} 给出读入设计并为顶层设计打开设计原理图后的 Design Vision 主窗口示例。

\figplaceholder{Figure 68: Design Vision Main Window}{Design Vision 主窗口}{fig:dv-main}

可在状态栏中显示菜单命令所执行操作的简要说明：将指针悬停在命令名上即可。
要确定工具栏按钮的功能，将指针悬停在按钮上；ToolTip 显示按钮名称，状态栏显示其用途的简要说明。

对于也可通过工具栏按钮或键盘快捷键使用的菜单命令，菜单会显示这些替代方式的示意。
部分常用菜单命令也出现在各视图的弹出菜单上；要使用弹出菜单，右键单击并选择命令。

可分别配置每个窗口的工具栏、视图窗口与面板。
工具栏始终附着于窗口边缘；可将个别面板附着于窗口边缘，或使其在工作区内/外浮动。
多数面板与特定类型的视图关联，并对活动视图操作；例外是控制台，它包含命令行及其自身视图。

可通过调整视图窗口与面板大小来改善工作环境；可最小化个别视图窗口，或最大化视图窗口以填满工作区；还可在工作区内平铺排列已打开的视图窗口。

要在下次会话中以相同配置恢复 GUI，必须在关闭 GUI 或退出 DC Explorer 工具之前保存当前窗口配置。
窗口配置包括主窗口与版图窗口的大小与位置，以及各工具栏与面板的可见性与位置。

关于 GUI 窗口的更多信息，见下列主题：
\begin{itemize}
  \item 输入命令并查看结果
  \item 查看 Man 页
  \item 保存窗口或视图的图像
  \item 在 GUI 中获取帮助
\end{itemize}

% ------------------------------------------------------------
\subsection{输入命令并查看结果}
\subsubsection*{Entering Commands and Viewing Results}

打开 GUI 时，控制台面板默认停靠在状态栏上方。可用控制台：
\begin{itemize}
  \item 在控制台命令行输入 \cmd{de\_shell} 命令
  \item 点击命令行上方的选项卡，查看 session transcript（log view）或 command history list（history view）
  \item 复制、编辑或重用命令，并搜索、选择与保存命令或消息
\end{itemize}

可在控制台命令行输入任意 \cmd{de\_shell} 命令。
发出命令后（按 Return，或点击命令行左侧的 prompt 按钮），console 与 \cmd{de\_shell} 都会回显命令输出，包括处理消息以及所有 warning 或 error message。
例如，若输入 \cmd{get\_selection} 命令，控制台 log 视图会显示所有所选对象的名称列表。

可用方向键在控制台命令行上显示、编辑并重新发出命令：左右移动插入点，或上下滚动命令栈。

log 视图显示会话信息转录，包括自打开 GUI 以来运行的命令及其输出与消息。可用这些信息：
\begin{itemize}
  \item 在执行功能后检查工具状态
  \item 排查遇到的问题
  \item 在转录中搜索以往功能相关信息
  \item 将转录、所选文本或仅错误与警告消息保存为文本文件
\end{itemize}

可像在 Linux shell 中一样，在 log 视图中复制文本并粘贴到命令行：用鼠标左键选择文本，用中键粘贴。

history 视图列出当前 GUI 会话中使用过的菜单、对话框与 \cmd{de\_shell} 命令。可对该列表：
\begin{itemize}
  \item 查看已使用的命令
  \item 查找并重新发出已使用的命令
  \item 保存列表供以后使用
\end{itemize}

可在 history 视图中选择命令，并在命令行上编辑或重新发出。

更多信息亦见 \textit{Design Vision User Guide} 中下列主题：
\begin{itemize}
  \item Viewing the Session Log
  \item Viewing the Command History
  \item Entering Commands in the Console
\end{itemize}

% ------------------------------------------------------------
\subsection{查看 Man 页}
\subsubsection*{Viewing Man Pages}

GUI 提供基于 HTML 的浏览器窗口，用于查看命令、变量与错误消息的 man 页。可用它：
\begin{itemize}
  \item 显示 man 页
  \item 在当前查看的 man 页中搜索文本
  \item 打印当前查看的 man 页
  \item 在已查看的 man 页之间前后浏览
\end{itemize}

要打开 man 页查看器：
\begin{itemize}
  \item 选择 \textbf{Help $>$ Man Pages}。
\end{itemize}
主页显示不同 man 页类别（commands、variables、messages）的链接列表。
点击要查看的类别链接，再点击要查看的 man 页标题链接。

也可在控制台命令行使用 \cmd{man} 或 \cmd{gui\_show\_man\_page} 命令，或在 GUI 已打开时在 \cmd{de\_shell} 中使用 \cmd{gui\_show\_man\_page} 命令，在 man 页查看器中显示 man 页。

更多信息亦见 \textit{Design Vision User Guide} 中的 ``Viewing Man Pages'' 主题。

% ------------------------------------------------------------
\subsection{保存窗口或视图的图像}
\subsubsection*{Saving an Image of a Window or View}

可将 GUI 窗口或视图窗口的图像保存为图像文件。
图像显示屏幕上所见窗口外观，但不含窗口边框或标题栏。
例如，若保存活动原理图视图的图像，图像显示当前缩放级别与平移位置下原理图的可见部分。

要保存活动 GUI 窗口或视图窗口的图像：
\begin{enumerate}
  \item 选择 \textbf{View $>$ Save Screenshot As}。\\
        出现 Save Screenshot As 对话框。
  \item 指定文件名与图像格式。\\
        图像格式可为 PNG（默认）、BMP、JPEG 或 XPM。要指定非默认格式，使用该格式的文件扩展名。
  \item （可选）若要保存活动视图窗口而非当前工作的 GUI 窗口的图像，选中 ``Grab screenshot of active view only'' 选项。
  \item 点击 \textbf{Save}。
\end{enumerate}

也可使用 \cmd{gui\_write\_window\_image} 命令，指定文件名、图像格式与窗口名。
窗口实例名出现在窗口标题栏与 Window 菜单上。例如，要将活动原理图视图保存为名为 \texttt{schem\_1.jpg} 的 JPEG 图像，输入：
\begin{lstlisting}
de_shell> gui_write_window_image -file schem_1.jpg
\end{lstlisting}

不能保存对话框或其他 GUI 元素（如工具栏或面板）的图像。

更多信息亦见 \textit{Design Vision User Guide} 中的 ``Saving an Image of a Window'' 主题。

% ------------------------------------------------------------
\subsection{在 GUI 中获取帮助}
\subsubsection*{Getting Help in the GUI}

Design Vision Help 可在 DC Explorer GUI 中使用。
关于 GUI 的多数所需信息可在面向 DC Explorer 的 Design Vision Help 中找到。
可从 Design Vision 主窗口或版图窗口的 Help 菜单访问 Design Vision Help。

要查看面向 DC Explorer 的 Design Vision Help：
\begin{itemize}
  \item 选择 \textbf{Help $>$ Online Help}。
\end{itemize}

要查看关于版图窗口中交互工具用法的 Design Vision Help 提示：
\begin{itemize}
  \item 选择 \textbf{Help $>$ Layout Help}。
\end{itemize}

帮助系统包含解释可执行任务细节的主题。例如，若用户指南中某步骤需要帮助，可在 Design Vision Help 中找到所需信息。

Design Vision Help 中的信息按下列类别组织：
\begin{itemize}
  \item Feature 主题：概述 Design Vision 窗口组件与工具
  \item How-to 主题：提供完成综合与分析任务的步骤
  \item Reference 主题：解释视图、工具栏按钮、菜单命令与对话框选项
\end{itemize}

Design Vision Help 是基于浏览器的 HTML Help 系统，设计用于在 Firefox 与 Mozilla Web 浏览器中查看。

\begin{noteBox}
在 Design Vision 内访问 Design Vision Help 之前，Web 浏览器可执行文件必须列在 Linux 的 \texttt{\$PATH} 变量中。

Design Vision Help 大量使用 JavaScript 与级联样式表（CSS）。若浏览器显示 Design Vision Help 时遇到问题，请打开浏览器首选项，确保已启用 JavaScript 与样式表，且安全首选项未阻止 JavaScript。
\end{noteBox}

关于 GUI 的其他信息来源包括 \cmd{gui\_*} 命令 man 页与 SolvNetPlus 知识库。

更多信息亦见 \textit{Design Vision User Guide} 中的 ``Using This Help System'' 主题。

% ============================================================
\section{执行 Floorplan 探索}
\subsection*{Performing Floorplan Exploration}

可在综合环境内，通过 IC Compiler 版图窗口中的 IC Compiler 设计规划工具，执行 floorplan 分析、创建与修改等 floorplanning 任务。
DC Explorer 中的 floorplan 探索与 IC Compiler 版图窗口之间的接口是透明的，可在 Design Vision 与 IC Compiler 版图窗口之间无缝切换。

要了解如何执行 floorplanning 任务，见：
\begin{itemize}
  \item 启用 Floorplan 探索
  \item 启用数据流分析
  \item 使用 Floorplan 探索 GUI
  \item 保存 Floorplan 或丢弃更新
  \item 退出会话
  \item 在 Floorplan 变更后执行综合
  \item 在批处理模式下使用 Floorplan 探索
  \item 黑盒、物理层次与块抽象
\end{itemize}

% ------------------------------------------------------------
\subsection{启用 Floorplan 探索}
\subsubsection*{Enabling Floorplan Exploration}

DC Explorer 允许您使用 IC Compiler 版图窗口中的 IC Compiler 设计规划工具执行 floorplan 探索。
DC Explorer 中的 floorplan 探索与 IC Compiler 版图窗口之间的接口是透明的。可在下列设计阶段执行 floorplan 探索：
\begin{itemize}
  \item \textbf{综合之前}\\
        要执行 floorplan 探索，确保变量 \varname{enable\_presynthesis\_floorplanning} 设为 \texttt{true}（默认）。
\begin{noteBox}
必须有逻辑库以及至少包含 AND 门、反相器与缓冲器的物理库，才能在 IC Compiler 中启动设计规划。
\end{noteBox}
  \item \textbf{综合之后}\\
        在执行 floorplan 探索之前，先在 DC Explorer 中读入 RTL 文件、指定物理与逻辑库、定义 Synopsys 设计约束、指定 floorplan（若有），并综合设计。
        综合后，可评估结果质量，并按需使用 floorplan 探索创建新 floorplan 或改进现有 floorplan。
\end{itemize}

要设置并启动 floorplan 探索会话：
\begin{enumerate}
  \item （可选）若要在综合前调用 floorplan 探索，将变量 \varname{enable\_presynthesis\_floorplanning} 设为 \texttt{true}。
  \item 在 DC Explorer GUI 中，选择 \textbf{Floorplan $>$ Set Design Planning Options}。\\
        出现 Set Design Planning Options 对话框，如图~\ref{fig:set-dp-opts} 所示。
  \item 在 Work Directory 框中，输入工作目录名及其相对于当前目录的路径。\\
        此处存放 floorplanning 脚本、floorplan 及其他必要文件。
  \item 在 ICC Executable 框中，输入 IC Compiler 可执行文件的名称与位置。\\
        默认情况下，DC Explorer 使用 \texttt{\$path} 变量指定的 \cmd{icc\_shell} 可执行文件。
  \item 在 File Name Prefix 框中，输入要用于任何生成文件（如 floorplanning 脚本与 floorplan 文件）文件名的前缀。\\
        默认文件前缀命名风格为 \texttt{\%s\_\%d}，其中 \texttt{\%s} 为设计名，\texttt{\%d} 为进程 ID。
  \item 点击 \textbf{OK}。
  \item 选择 \textbf{Floorplan $>$ Start Design Planning}。\\
        出现警告，显示下列消息：
\begin{lstlisting}
Warning: Starting floorplan exploration will close the GUI. Do you
 want to continue?
\end{lstlisting}
  \item 点击 \textbf{OK}。\\
        DC Explorer 在 IC Compiler layout window 中启动 floorplan exploration，并使用简化的 floorplanning menu。
\end{enumerate}

\figplaceholder{Figure 69: Set Design Planning Options Dialog Box}{Set Design Planning Options 对话框}{fig:set-dp-opts}

作为步骤 1 至 8 的替代，可使用下列 \cmd{de\_shell} 命令：
\begin{lstlisting}
de_shell> set_icc_dp_options -work_dir "./dcg_link" \
   -icc_executable "*/bin/icc_shell" \
   -file_name_prefix "%s_%d" -keep_files
de_shell> start_icc_dp
\end{lstlisting}

更多信息亦见：
\begin{itemize}
  \item \textit{Design Vision User Guide} 中的 ``Using Floorplan Exploration Tools'' 主题
  \item IC Compiler Help
\end{itemize}

\begin{seeAlsoBox}
在批处理模式下使用 Floorplan 探索
\end{seeAlsoBox}

% ------------------------------------------------------------
\subsection{启用数据流分析}
\subsubsection*{Enabling Data Flow Analysis}

Floorplan 探索允许您使用 IC Compiler 设计规划中的数据流分析工具，分析并改进物理设计中的数据流。
可用数据流分析器窗口执行下列数据流分析：
\begin{itemize}
  \item logic connectivity analysis：从 logic data flow 视角查看 macro 连接
  \item 高级飞线分析（advanced flyline analysis）：通过检查宏之间以及宏与 I/O 单元之间的连接来识别数据流
  \item 交互式宏编辑：修改宏布局、创建宏组与宏阵列
\end{itemize}

图~\ref{fig:early-fp-dfa} 给出从 DC Explorer 到 IC Compiler 设计规划的早期 floorplanning 概览。

\figplaceholder{Figure 70: Early Floorplanning Using Data Flow Analysis}{使用数据流分析的早期 Floorplanning}{fig:early-fp-dfa}

要启动数据流分析会话：
\begin{enumerate}
  \item 按「启用 Floorplan 探索」中的步骤调用 IC Compiler 版图窗口。\\
        出现 IC Compiler 版图窗口。
  \item 选择 \textbf{Floorplan $>$ Analyze Logic Connectivity}。\\
        出现数据流分析器窗口。
  \item 使用 IC Compiler floorplanning 工具执行 floorplanning 任务，包括创建、修改与分析 floorplan。
  \item 选择 \textbf{Floorplan $>$ Update DC Floorplan}，将 floorplanning 更新保存回 DC Explorer。
  \item （可选）按「使用 Floorplan 物理约束执行优化」中的步骤执行优化。
\end{enumerate}

更多信息亦见：
\begin{itemize}
  \item IC Compiler Help 中的 ``Analyzing the Data Flow for Macro Placement'' 主题
  \item \textit{IC Compiler Implementation User Guide}
\end{itemize}

\begin{seeAlsoBox}
\begin{itemize}
  \item 启用 Floorplan 探索
  \item 使用 Floorplan 物理约束执行优化
  \item SolvNetPlus 文章 24224，``Design Compiler Graphical Floorplan Exploration Application Note''
\end{itemize}
\end{seeAlsoBox}

% ------------------------------------------------------------
\subsection{使用 Floorplan 探索 GUI}
\subsubsection*{Using the Floorplan Exploration GUI}

DC Explorer 中的 floorplan 探索使用 IC Compiler 版图窗口。
默认情况下，启用 floorplan 探索时，DC Explorer 会配置 IC Compiler 版图窗口，使其仅包含执行 DC Explorer floorplanning 所需的菜单与菜单命令。
图~\ref{fig:icc-simple-layout} 给出简化 floorplanning 菜单结构。

\figplaceholder{Figure 71: IC Compiler Simplified Layout Window Features}{IC Compiler 简化版图窗口功能}{fig:icc-simple-layout}

若您是有经验的 IC Compiler 用户，可使用完整的设计规划菜单集。

可通过选择 \textbf{File $>$ Task $>$ Design Planning} 访问完整的 IC Compiler 设计规划菜单集，如图~\ref{fig:switch-dp-menus} 所示。
要从完整设计规划菜单结构切换回简化菜单结构，选择 \textbf{File $>$ Task $>$ DCG-DP Link}。

\figplaceholder{Figure 72: Switching Between Simplified and Full Design Planning Menu Structures}{在简化与完整设计规划菜单结构之间切换}{fig:switch-dp-menus}

使用完整 IC Compiler 设计规划菜单集完成的任何工作，在切换回简化 floorplanning 菜单后仍然可用。

\subsubsection{创建与编辑 Floorplan}
\paragraph*{Creating and Editing Floorplans}

无论使用简化的 DC Explorer floorplanning 菜单还是完整的 IC Compiler 设计规划菜单，IC Compiler 版图窗口都提供可用于创建或修改 floorplan 的交互工具与命令。您可以：
\begin{itemize}
  \item 使用编辑工具移动、调整大小、复制、对齐、分布、展开、分割、重塑或移除对象
  \item 使用编辑命令旋转、对齐、分布、展开或扩展对象，或更改单元方向
\end{itemize}

最常用的编辑操作提供快捷键；每个编辑工具都提供联机帮助页。

此外，可用 Floorplan 菜单上的命令创建物理对象，例如 placement blockage、bound 与 routing keepout。
创建对象时，可指定其坐标，或在版图视图中绘制该对象。

更多信息亦见：
\begin{itemize}
  \item IC Compiler Help
  \item \textit{IC Compiler Design Planning User Guide}
\end{itemize}

% ------------------------------------------------------------
\subsection{保存 Floorplan 或丢弃更新}
\subsubsection*{Saving the Floorplan or Discarding Updates}

可随时从 IC Compiler floorplanning 会话中保存或丢弃 DC Explorer 中的 floorplan 更新。要保存或丢弃 floorplan，执行下列步骤：
\begin{enumerate}
  \item 选择 \textbf{Floorplan $>$ Update DC Floorplan}。\\
        出现 Update DC Floorplan 对话框，如图~\ref{fig:update-dc-fp} 所示。
\begin{noteBox}
``Update DC Floorplan'' 命令仅在面向 DC Explorer 的简化 floorplanning 菜单结构中可用。
关于从完整设计规划菜单集切换到简化 floorplanning 菜单结构的信息，见「使用 Floorplan 探索 GUI」。
\end{noteBox}
  \item 选择适当选项。
  \item 点击 \textbf{OK}。
\end{enumerate}

\figplaceholder{Figure 73: Update DC Floorplan Dialog Box}{Update DC Floorplan 对话框}{fig:update-dc-fp}

也可将 floorplan 保存在下列文件中：
\begin{itemize}
  \item Tcl 脚本文件\\
        选择 \textbf{Floorplan $>$ Write Floorplan}。也可在工具提示符下使用 \cmd{write\_floorplan} 命令。
  \item DEF 文件\\
        选择 \textbf{File $>$ Export $>$ Write DEF}。也可在工具提示符下使用 \cmd{write\_def} 命令。
\end{itemize}

% ------------------------------------------------------------
\subsection{退出会话}
\subsubsection*{Exiting the Session}

要退出 floorplan 探索会话并保存 floorplan 供综合使用，或丢弃 floorplan，执行下列步骤：
\begin{enumerate}
  \item 选择 \textbf{File $>$ Exit}。\\
        出现 Exit IC Compiler 对话框，如图~\ref{fig:exit-icc} 所示。
  \item 选择适当选项。
  \item 点击 \textbf{OK}。
\end{enumerate}

\figplaceholder{Figure 74: Exit IC Compiler Dialog Box}{Exit IC Compiler 对话框}{fig:exit-icc}

若更新了 floorplan，工具会关闭原来的 Design Vision 版图窗口。要查看 floorplan 变更，请打开新的 Design Vision 版图窗口。

% ------------------------------------------------------------
\subsection{在 Floorplan 变更后执行综合}
\subsubsection*{Performing Synthesis After Floorplan Changes}

要在 floorplan 变更后准备完整综合：
\begin{enumerate}
  \item 必须运行 \cmd{write\_floorplan} 命令，将 floorplan 保存为 Tcl 脚本文件。例如：
\begin{lstlisting}
de_shell> write_floorplan -all Design.fp
\end{lstlisting}
        对使用新的或已修改的 floorplan 重新综合设计，建议使用 \opt{-all} 选项。
  \item 在 DC Explorer 中用 \cmd{read\_floorplan} 命令读入 floorplan 以供综合。
  \item 读入 RTL 文件、SDC 文件与新的 floorplan 文件后，用 \cmd{compile\_exploration} 命令执行完整综合。
\end{enumerate}

\begin{seeAlsoBox}
使用 Floorplan 物理约束
\end{seeAlsoBox}

% ------------------------------------------------------------
\subsection{在批处理模式下使用 Floorplan 探索}
\subsubsection*{Using Floorplan Exploration in Batch Mode}

也可在批处理模式下使用 floorplan 探索。批处理模式允许您在命令行用脚本创建并修改 floorplan。

在 DC Explorer 中用下列命令以批处理模式执行 floorplan 探索：
\begin{enumerate}
  \item 使用 \cmd{set\_icc\_dp\_options} 命令指定 IC Compiler 可执行文件及其他设置要求。例如：
\begin{lstlisting}
de_shell> set_icc_dp_options -icc_executable icc_image
\end{lstlisting}
  \item （可选）使用 \cmd{report\_icc\_dp\_options} 命令报告 \cmd{set\_icc\_dp\_options} 所指定的选项。\\
        若未对 \cmd{set\_icc\_dp\_options} 指定任何选项，\cmd{report\_icc\_dp\_options} 报告默认设置。
  \item 使用 \cmd{start\_icc\_dp} 命令启动 floorplan 探索会话。\\
        \cmd{start\_icc\_dp} 使用 \cmd{set\_icc\_dp\_options} 指定的 IC Compiler 可执行文件。
        要在批处理模式下 source 脚本，对 \cmd{start\_icc\_dp} 使用 \opt{-f} 选项。例如：
\begin{lstlisting}
de_shell> start_icc_dp -f ../../scripts/script_file_name.tcl
\end{lstlisting}
\end{enumerate}

\begin{seeAlsoBox}
启用 Floorplan 探索
\end{seeAlsoBox}

% ------------------------------------------------------------
\subsection{黑盒、物理层次与块抽象}
\subsubsection*{Black Boxes, Physical Hierarchies, and Block Abstractions}

DC Explorer 支持对包含黑盒（black box）、物理层次（physical hierarchy）与块抽象（block abstraction）的网表进行 floorplan 探索。
\begin{itemize}
  \item \textbf{黑盒}\\
        若设计含有黑盒，无需执行额外步骤或修改 floorplan 探索流程。
        在 DC Explorer 中定义或创建的黑盒会自动传输到 IC Compiler 设计规划。
        使用 IC Compiler 工具，可在 floorplanning 会话中编辑含黑盒的 floorplan，并将修改传回 DC Explorer。
  \item \textbf{物理层次与块抽象}\\
        在 floorplan 探索期间，DC Explorer 会自动将物理层次与块抽象传输到 IC Compiler 设计规划。
        物理层次与 block abstraction 作为 macro 传输，block 内包含 quick timing model。\\
        在将物理层次与块抽象传输到 IC Compiler 工具以进行设计规划时，DC Explorer 会自动在实例上将 \texttt{dct\_physical\_hier} 属性设为 \texttt{true}。
\end{itemize}

\subsubsection{查询层次块与引脚连接}
\paragraph*{Querying Hierarchical Blocks and Pin Connections}

要返回所有 DC Explorer 层次块（包括物理层次与块抽象），对 \cmd{get\_cells} 命令使用 \texttt{dct\_physical\_hier} 属性。
将物理层次与块抽象传回 DC Explorer 时会予以保留。

例如：
\begin{lstlisting}
icc_shell> get_cells -hierarchical "dct_physical_hier==true"
{GPRs Multiplier}
\end{lstlisting}

要查询连接到层次块引脚的网，对 \cmd{get\_nets} 使用 \opt{-boundary\_type lower} 或 \opt{-boundary\_type upper} 选项，分别针对层次块内部或外部的网。
要同时查询两者，指定 \texttt{both} 关键字。
若不指定 \opt{-boundary\_type} 选项，命令返回层次块外部的网。
使用该选项时，必须用 \opt{-of\_objects} 选项指定引脚。

例如，下列两条命令返回连接到 \texttt{u2} 层次块外部 \texttt{in} 引脚的 \texttt{topnet} 网：
\begin{lstlisting}
de_shell> get_nets -of_objects u2/in
{topnet}
de_shell> get_nets -of_objects u2/in -boundary_type upper
{topnet}
\end{lstlisting}

下列命令返回连接到 \texttt{u2} 层次块内部 \texttt{in} 引脚的 \texttt{u2/in} 网：
\begin{lstlisting}
de_shell> get_nets -of_objects u2/in -boundary_type lower
{u2/in}
\end{lstlisting}

% ============================================================
\section{使用 IC Compiler II 进行 Floorplan 探索}
\subsection*{Floorplan Exploration With IC Compiler II}

DC Explorer 允许您使用 IC Compiler II 版图窗口中的 IC Compiler II 设计规划工具执行 floorplan 探索。
综合后，可评估结果质量，并在 IC Compiler II 中使用 floorplan 探索创建 floorplan 或改进现有 floorplan。
DC Explorer 中的 floorplan 探索与 IC Compiler II 版图窗口之间的接口是透明的。

下列主题说明如何用 IC Compiler II 执行 floorplan 探索：
\begin{itemize}
  \item Floorplan 探索的先决条件
  \item 运行 Floorplan 探索流程
  \item 向 IC Compiler II 传输数据
  \item 以 ASCII 格式保存以供 IC Compiler II 使用
\end{itemize}

\begin{seeAlsoBox}
执行 Floorplan 探索
\end{seeAlsoBox}

% ------------------------------------------------------------
\subsection{Floorplan 探索的先决条件}
\subsubsection*{Prerequisites for Floorplan Exploration}

在 IC Compiler II 中执行 floorplan 探索之前，必须完成下列任务：
\begin{itemize}
  \item \textbf{综合设计}\\
        不能对未综合（即 unmapped）的设计运行 IC Compiler II。
  \item \textbf{生成 IC Compiler II 参考库}\\
        由于参考库不会由 DC Explorer 自动转换，必须使用 IC Compiler II Library Manager 生成库。
  \item \textbf{为多电压设计生成 RTL UPF 文件}\\
        为处理设计的功耗意图，运行 RTL UPF 探索流程以生成 RTL UPF 文件。
\end{itemize}

\begin{seeAlsoBox}
UPF 探索
\end{seeAlsoBox}

% ------------------------------------------------------------
\subsection{运行 Floorplan 探索流程}
\subsubsection*{Running the Floorplan Exploration Flow}

下图给出从 DC Explorer 到 IC Compiler II 设计规划的 floorplan 探索概览：

\figplaceholder{Figure 75: From DC Explorer to Floorplan Exploration in IC Compiler II}{从 DC Explorer 到 IC Compiler II 中的 Floorplan 探索}{fig:dce-to-icc2-fp}

下列各节说明先决条件，以及以交互模式或以 \cmd{de\_shell} 脚本批处理模式运行的 floorplan 探索流程。

\subsubsection{先决条件}
\paragraph*{Prerequisites}

在 floorplan 探索之前，必须完成下列任务：
\begin{enumerate}
  \item 用 \cmd{compile\_exploration} 命令综合设计。
  \item 用 \cmd{set\_icc2\_options} 命令指定 IC Compiler II 参考库。例如：
\begin{lstlisting}
de_shell> set_icc2_options -icc2_executable ../icc2_shell \
   -ref_libs "lib1.ndm lib2.ndm"
\end{lstlisting}
\end{enumerate}

\subsubsection{交互模式}
\paragraph*{In Interactive Mode}

用 \cmd{start\_icc2} 命令运行 floorplan 探索流程需要 IC Compiler II 许可证。

要启动 floorplan 探索流程：
\begin{enumerate}
  \item 在 DC Explorer 中，用 \cmd{start\_icc2} 命令启动 IC Compiler II 设计规划会话。\\
        \cmd{start\_icc2} 调用 \texttt{icc2\_shell>} 提示符并打开 IC Compiler II 版图窗口。
  \item 在 IC Compiler II 中执行 floorplan 探索。\\
        可用 IC Compiler II 版图窗口中的设计规划工具创建 floorplan 或编辑现有 floorplan。
  \item 在 \cmd{icc2\_shell} 中用 \cmd{update\_dc\_floorplan} 命令生成更新后的 floorplan。\\
        IC Compiler II 中的 floorplan 数据会传输到 DC Explorer，并 reset placement。
        若不执行此步骤，在 \cmd{icc2\_shell} 中所做的 floorplan 变更不会更新到 \cmd{de\_shell}，并将丢失。
  \item （可选）要在步骤 3 中保留 floorplan 文件，对 \cmd{set\_icc2\_options} 指定 \opt{-keep\_files} 选项。
  \item 在 \cmd{icc2\_shell} 中输入 \cmd{exit} 命令以结束设计规划会话。\\
        返回 \texttt{de\_shell>} 提示符。
\end{enumerate}

\begin{noteBox}
运行 \cmd{start\_icc2} 命令时，若指定了 \texttt{TMPDIR} Linux 环境变量，则调用 IC Compiler II 所需的文件会写入 \texttt{\$TMPDIR/dcg\_unique\_string} 工作目录。
要更改默认目录，对 \cmd{set\_icc2\_options} 指定 \opt{-work\_dir} 选项。
\end{noteBox}

\subsubsection{批处理模式}
\paragraph*{In Batch Mode}

要从批处理脚本运行 floorplan 探索，在启动 DC Explorer 时对 \cmd{start\_icc2} 指定 \opt{-f script\_file} 选项。
在批处理模式下，\cmd{start\_icc2} 调用 \cmd{icc2\_shell} 并执行指定脚本，但不打开 GUI。

下列脚本给出与 IC Compiler II 配合的一般 floorplan 探索流程示例：
\begin{lstlisting}
#Perform synthesis using compile_exploration
compile_exploration -scan -gate_clock

#Set the IC Compiler II options
set_icc2_options -icc2_executable ../icc2_shell \
-ref_libs "lib1.ndm lib2.ndm"

#Launch IC Compiler II session in batch mode
start_icc2 -f my_icc2.tcl
\end{lstlisting}

示例批处理脚本 \texttt{my\_icc2.tcl}：
\begin{lstlisting}
#This example script displaces the macro I_RAM by 10 microns
set_placement_status placed [get_cells I_RAM]
move_objects -delta {10 10} [get_cells I_RAM]
set_placement_status fixed [get_cells I_RAM]

#Update the floorplan changes before exit
update_dc_floorplan
exit
\end{lstlisting}

% ------------------------------------------------------------
\subsection{向 IC Compiler II 传输数据}
\subsubsection*{Data Transfer to IC Compiler II}

在 DC Explorer 内启动 IC Compiler II 时，设计数据会以 IC Compiler II 兼容格式自动传输到 IC Compiler II。
表~\ref{tab:xfer-icc2} 列出从 DC Explorer 传输到 IC Compiler II 的设计与环境设置。

\begin{table}[htbp]
\centering
\caption{传输到 IC Compiler II 的设计与环境设置}
\label{tab:xfer-icc2}
\begin{tabularx}{\textwidth}{@{}l X@{}}
\toprule
\textbf{传输} & \textbf{设计与环境设置} \\
\midrule
已传输 &
\begin{itemize}[leftmargin=1.2em,nosep]
  \item Verilog 网表
  \item IC Compiler II 兼容格式的时序约束
  \item DEF 文件与 IC Compiler II 兼容格式 Tcl floorplan 文件中的 floorplan 约束
  \item 属性，例如 \texttt{dont\_touch}、\texttt{size\_only}、\texttt{dont\_use}
  \item Technology 文件
  \item TLUPlus 文件
  \item UPF 文件
  \item 与层相关的约束，例如网布线层约束与首选布线方向
  \item 忽略层（ignored layer）设置
\end{itemize} \\
\midrule
未传输 &
\begin{itemize}[leftmargin=1.2em,nosep]
  \item hierarchical model：DC Explorer 与 IC Compiler 的 block abstraction 和 physical hierarchy 在 \cmd{icc2\_shell} 中被视为 unresolved reference
  \item 物理黑盒的估算：这些物理单元由 \cmd{estimate\_fp\_black\_boxes} 命令创建
  \item 相对布局约束：这些约束在 DC Explorer 中创建
\end{itemize} \\
\bottomrule
\end{tabularx}
\end{table}

% ============================================================
\section{分析 RTL}
\subsection*{Analyzing the RTL}

可在早期设计阶段通过分析设计的整体时序与逻辑概况，来调试时序与拥塞问题。
DC Explorer GUI 使您能生成此类概况，以识别最差时序路径、逻辑级数最多的路径、不可行路径、多周期路径等。
要分析设计，可生成直方图以显示落在各 slack 范围或逻辑级范围（称为 bin）中的路径分布。
选定某个 bin 后，可将路径载入路径数据表，并在 RTL 浏览器中对其单元做交叉探测。

\begin{noteBox}
要使用此分析功能，设计必须已 mapped 或已 compile。
\end{noteBox}

在 DC Explorer GUI 中，可执行下列 RTL 分析：
\begin{itemize}
  \item \textbf{生成路径 slack 直方图}
  \begin{enumerate}
    \item 选择 \textbf{AnalyzeRTL $>$ Analyze Path Slack}。
    \item 选择 \textbf{Single Clock Paths}、\textbf{Cross Clock Paths}、\textbf{Infeasible Paths} 或 \textbf{Multicycle Paths}。
  \end{enumerate}
\end{itemize}

\figplaceholder{Figure 76: Path Slack Histogram}{路径 Slack 直方图}{fig:path-slack-hist}

可创建违及时序约束的同时钟路径、跨时钟路径、不可行路径或多周期路径的直方图。
默认情况下，path slack histogram 使用 10 个 bin 显示所有路径的 slack 值分布，每个 bin 的最小区间宽度为 100~ps。
各 bin 表示路径数量（$y$ 轴）相对于其 slack 值（$x$ 轴）。

选定某个 bin 后，可选择该 bin 中的路径、在新的路径 slack 或逻辑级直方图中显示这些路径，或将路径载入路径数据表。

\begin{itemize}
  \item \textbf{生成逻辑级直方图}
  \begin{enumerate}
    \item 选择 \textbf{AnalyzeRTL $>$ Analyze Logic Levels}。
    \item 选择 \textbf{Single Clock Paths}、\textbf{Cross Clock Paths}、\textbf{Infeasible Paths} 或 \textbf{Multicycle Paths}。
  \end{enumerate}
\end{itemize}

\figplaceholder{Figure 77: Logic-Level Histogram}{逻辑级直方图}{fig:logic-level-hist}

可创建违及时序约束的同时钟路径、跨时钟路径、不可行路径或多周期路径的直方图。
默认情况下，logic-level histogram 使用 8 至 16 个 bin 显示所有路径的 logic level 分布。
各 bin 表示路径数量（$y$ 轴）相对于逻辑级数（$x$ 轴）。

选定某个 bin 后，可选择该 bin 中的路径、在新的路径 slack 或逻辑级直方图中显示这些路径，或将路径载入路径数据表。

\begin{itemize}
  \item \textbf{生成路径数据表}
  \begin{enumerate}
    \item 在路径 slack 或逻辑级直方图中选择一个或多个 bin。\\
          所选 bin 的颜色变为黄色。
    \item 右键单击并选择 \textbf{Create Data Table}。
  \end{enumerate}
\end{itemize}

\figplaceholder{Figure 78: Logic-Level Histogram and Path Data Table}{逻辑级直方图与路径数据表}{fig:logic-level-table}

\begin{itemize}
  \item \textbf{对单元做交叉探测}
  \begin{enumerate}
    \item 选择一个或多个单元、引脚或端口。\\
          可在原理图或版图视图、层次浏览器、对象列表中选择对象，或通过 Select 菜单上的命令选择。
    \item 用下列任一方式做交叉探测：
    \begin{itemize}
      \item 选择 \textbf{AnalyzeRTL $>$ Cross Probe to Source}。
      \item 右键单击所选对象并选择 \textbf{Cross Probe to Source}。
    \end{itemize}
  \end{enumerate}
        工具定位单元来源的 RTL 文件，并在 RTL 浏览器中打开这些文件。
\end{itemize}

\begin{itemize}
  \item \textbf{对时序路径做交叉探测}
  \begin{enumerate}
    \item 在路径数据表中选择一条或多条路径。
    \item 用下列任一方式做交叉探测：
    \begin{itemize}
      \item 选择 \textbf{AnalyzeRTL $>$ Cross Probe to Source}。
      \item 右键单击所选路径并选择 \textbf{Cross Probe to Source}。
    \end{itemize}
  \end{enumerate}
        工具为每条路径上的所有单元创建集合，定位单元来源的 RTL 文件，并在新的 RTL 浏览器窗口中打开这些文件。
\end{itemize}

\figplaceholder{Figure 79: Cross-Probing Between Path Data Table and RTL Browser}{路径数据表与 RTL 浏览器之间的交叉探测}{fig:xprobe-path-rtl}

可在下列各设计视图与 RTL 浏览器之间做交叉探测：
\begin{itemize}
  \item 原理图视图
  \item 路径 slack 或逻辑级直方图（经路径数据表）
  \item 路径检查器（path inspector）
  \item 拥塞图（congestion map）
\end{itemize}

更多信息亦见 \textit{Design Vision User Guide} 中的 ``Analyzing RTL'' 主题。

% ============================================================
\section{RTL 交叉探测}
\subsection*{RTL Cross-Probing}

DC Explorer 在 GUI 中提供 cross-probing 功能，使您能在 RTL browser 中审阅并验证设计对象与 RTL 源码之间的对应关系。
可从某一视图（如原理图、单元列表、版图或时序视图）交叉探测设计中的单元、引脚或时序路径，并检查对应的 RTL 文件。
这使您能识别可修改以改善时序、拥塞或 datapath 提取的代码行。

\begin{noteBox}
要使用交叉探测能力，必须有已 elaborate 或已综合的设计。
\end{noteBox}

可从路径检查器选择关键路径，工具会交叉探测 RTL 以帮助调试最差情况时序路径。
拥塞图与 RTL 浏览器之间的交叉探测有助于识别导致拥塞的 RTL 代码。
也可在路径数据表与 RTL 浏览器之间做交叉探测；路径数据表显示您在路径 slack 或逻辑级直方图中所检查时序路径的详细信息。

工具不支持下列交叉探测能力：
\begin{itemize}
  \item 从各种设计视图到 UPF 文件
  \item 选择网作为交叉探测对象
\end{itemize}

图~\ref{fig:xprobe-layout-rtl} 给出 Design Vision 版图视图与 RTL 之间的交叉探测：

\figplaceholder{Figure 80: Cross-Probing Between the Layout View and the RTL Browser}{版图视图与 RTL 浏览器之间的交叉探测}{fig:xprobe-layout-rtl}

% ------------------------------------------------------------
\subsection{重复 RTL 交叉探测}
\subsubsection*{Repeating RTL Cross-Probing}

每次工具为交叉探测读取 RTL 时，会将 RTL 行号与文件路径存入数据库。
若在 cross-probing 后将 RTL 文件移到其他位置，数据库中保存的文件路径将失效；再次执行同一 cross-probing 序列时，工具无法定位该文件并会发出错误消息。

要重复交叉探测序列或引用先前的交叉探测数据，必须先查明曾被交叉探测的文件状态。可通过运行下列命令获取状态：
\begin{itemize}
  \item \cmd{report\_cross\_probing\_files}\\
        该命令报告已交叉探测文件的状态，包括 OK、Missing、Modified 与 No Permissions。
  \item \cmd{update\_cross\_probing\_files}\\
        该命令更新曾被交叉探测的文件位置，并列出文件的先前路径与当前路径。
        若文件已被修改，该命令不会更新已交叉探测的文件。
\end{itemize}

更多信息亦见 \textit{Design Vision User Guide} 中下列主题：
\begin{itemize}
  \item Analyzing RTL
  \item Viewing Path Data Tables
\end{itemize}

\begin{seeAlsoBox}
\begin{itemize}
  \item 分析 RTL
  \item 使用命令进行 RTL 交叉探测
\end{itemize}
\end{seeAlsoBox}
